\pageoriginale

\clearpage
\SGAExposeBookmark{XXI}{El nivel de la cohomología de las intersecciones completas}
\begin{center}
{\large SGA 7 --- EXPOSÉ XXI}\par
\vspace{2em}
{\Large\underline{El nivel de la cohomología de las intersecciones
completas}}\par
\vspace{1.5em}
{\large por N. Katz}\par
\vspace{0.6em}
con un apéndice de P. Deligne (\S~5)
\end{center}

\section*{Sumario}
\begin{description}
\item[0.] Introducción \hfill 1
\item[1.] La matriz de Hasse--Witt de una intersección completa \hfill 2
\item[2.] El conivel sobre un cuerpo finito de definición \hfill 14
\item[3.] Aplicación a las intersecciones completas \hfill 16
\item[4.] Conclusiones \hfill 19
\item[5.] Teoremas de integralidad \hfill 21
\end{description}

\pageoriginale

\section*{0. Introducción}

Sea $X$ una variedad proyectiva lisa de dimensión $n\geq3$, sea $Y(d)$
una sección por una hipersuperficie de grado $d$, y sea
$E^{n-1}(Y(d),\mathbb Q_\ell)$ la cohomología «evanescente» de $Y(d)$
(XVIII~5.4.4). El objeto de esta exposición es demostrar que, para $d$
suficientemente grande e $Y(d)$ suficientemente general, no es cierto que
todo $E^{n-1}(Y(d),\mathbb Q_\ell)$ «provenga» de los grupos $H^i$,
$i<n-1$, de variedades algebraicas lisas y propias mediante
correspondencias algebraicas. En característica cero esto se deduce
inmediatamente de la bigraduación de Hodge de la cohomología compleja y
del hecho de que, en cohomología coherente,
\begin{equation*}
\dim H^{n-1}(Y(d),\mathcal O_{Y(d)})\longrightarrow\infty
\qquad\text{cuando }d\longrightarrow\infty.
\end{equation*}

En característica $p$, el método relaciona de nuevo el grupo coherente
$H^{n-1}(Y(d),\mathcal O_{Y(d)})$ con el grupo étale
$H^{n-1}_{\mathrm{et}}(Y(d),\mathbb Q_\ell)$. Para ello se trabaja sobre
un cuerpo base finito,

\pageoriginale

donde ambos quedan ligados por la función zeta de $Y(d)$: esta función se
expresa en términos de la cohomología $\ell$-ádica, mientras que su
«reducción módulo $p$» se expresa en términos de la cohomología coherente
(cf.~XXII~3.9). Para las intersecciones completas este método funciona sin
dificultad, gracias al conocimiento detallado de ambos tipos de
cohomología, étale y coherente, en este caso. Da el resultado principal,
el Teorema~4.2, como consecuencia de un resultado sobre las operaciones de
Hasse--Witt, el Teorema~1.4 (véase 1.0).

Hay argumentos heurísticos, debidos a A. Grothendieck y basados en el yoga
de la cohomología cristalina, que hacen plausible el enunciado general
para todo $X$ proyectivo liso, esencialmente por el mismo método.

Incluso en el caso particular de las intersecciones completas, la
demostración utiliza de manera esencial en 2.2 el Teorema~2.1 sobre la
integralidad de los valores propios de Frobenius, demostrado recientemente
por P. Deligne y demostrado aquí en el apéndice (\S~5).

\section*{1. La matriz de Hasse--Witt de una intersección completa}

\subsubsection*{Definición 1.0.}
Un esquema propio sobre un cuerpo $k$ de exponente característico $p$ se
llama \emph{especial} si la operación aditiva $p$-lineal
\begin{equation}\tag{1.0.1}
\mathfrak F_p:\mathcal O_X\longrightarrow\mathcal O_X,
\qquad g\longmapsto g^p,
\end{equation}
induce una operación nilpotente
\begin{equation}\tag{1.0.2}
\mathfrak F_p:H^{\dim X}(X,\mathcal O_X)
\longrightarrow H^{\dim X}(X,\mathcal O_X).
\end{equation}

\pageoriginale

Equivalentemente, si
$M=\dim_k H^{\dim X}(X,\mathcal O_X)$, la operación $p^M$-lineal
\begin{equation}\tag{1.0.3}
\mathfrak F_p^M=\mathfrak F_{p^M}:\mathcal O_X\longrightarrow\mathcal O_X,
\qquad g\longmapsto g^{p^M},
\end{equation}
anula $H^{\dim X}(X,\mathcal O_X)$.

Recordemos [2] que, para todo $i\in\mathbb Z$, el homomorfismo
\begin{equation*}
H^i(X,\mathcal O_X)\xrightarrow{\ \mathfrak F_p\ }
H^i(X,\mathcal O_X)
\end{equation*}
inducido por (1.0.1) se llama clásicamente \emph{homomorfismo de
Hasse--Witt} sobre $H^i(X,\mathcal O_X)$. Obsérvese que, cuando $p=1$,
decir que $X$ no es especial significa simplemente que
$H^{\dim X}(X,\mathcal O_X)\ne0$.

\subsubsection*{Teorema 1.1.}
Fijemos un exponente característico $p$, una dimensión $n\geq1$, un
entero $r\geq1$ y un multigrado $(d_1,\dots,d_r)$. Si
\begin{equation}\tag{1.1.0}
\sum_{i=1}^{r}d_i\geq(n+r)+1,
\end{equation}
entonces, entre todas las intersecciones completas de dimensión $n$ en
$\mathbb P^{n+r}$, de multigrado $(d_1,\dots,d_r)$ y característica $p$,
una suficientemente general no es especial en el sentido de 1.0.

\noindent\underline{Demostración.}
Sea
\begin{equation}\tag{1.1.1}
f:X_U\longrightarrow U
\end{equation}
la «familia universal» de intersecciones completas de multigrado
$(d_1,\dots,d_r)$ en $\mathbb P^{n+r}$ y característica $p$, provistas de
sus ecuaciones de definición. Así, $U$ es un abierto de un producto de
$r$ espacios proyectivos sobre $\mathbb F_p$, de dimensión impresionante.

\pageoriginale

El $\mathcal O_U$-módulo
\begin{equation}\tag{1.1.2}
R^nf_*\mathcal O_{X_U}
\end{equation}
es localmente libre y su formación conmuta con todo cambio de base
(XI~1.5(i)). Por (1.1.0), su rango $M$ no es nulo (XI~2.5). Si $p=1$,
esto demuestra el teorema para toda intersección completa que verifica
(1.1.0).

Supongamos ahora $p\geq2$. La operación iterada de Hasse--Witt
\begin{equation}\tag{1.1.3}
\mathfrak F_p^M=\mathfrak F_{p^M}:
R^nf_*\mathcal O_{X_U}\longrightarrow R^nf_*\mathcal O_{X_U}
\end{equation}
es $p^M$-lineal sobre $\mathcal O_U$. Como
$R^nf_*\mathcal O_{X_U}$ es localmente libre, (1.1.3) se representa
localmente sobre $U$ por una matriz $(M,M)$, denotada $H$, con coeficientes
en $\mathcal O_U$. Por consiguiente,
\begin{equation}\tag{1.1.4}
\{u\in U\mid X_u\text{ no es especial}\}\quad\text{es abierto en }U.
\end{equation}
En efecto, es el conjunto de los $u$ para los que $H(u)\ne0$. Como $U$ es
irreducible, basta hallar un valor de $u$ para el que $X_u$ no sea
especial; lo haremos sobre un cuerpo finito.

\subsubsection*{Proposición 1.1.5.}
Sea $X$ una intersección completa sobre $\mathbb F_q$. Si $X$ no tiene
puntos racionales sobre $\mathbb F_q$, entonces $X$ no es especial.

\pageoriginale

\noindent\underline{Demostración.}
Basta evidentemente demostrar, para una intersección completa $X$, la
implicación
\begin{equation}\tag{1.1.5.1}
X\text{ especial}\quad\Longrightarrow\quad
\operatorname{card}\bigl(X(\mathbb F_q)\bigr)\equiv1\pmod p.
\end{equation}
Por la fórmula de las trazas, demostrada en el Exposé XXII, 3.2,
\begin{equation*}
\operatorname{card}\bigl(X(\mathbb F_q)\bigr)
\equiv\sum_{i\geq0}(-1)^i
\operatorname{tr}\bigl(\mathfrak F_q\mid H^i(X,\mathcal O_X)\bigr)
\pmod p.
\end{equation*}
Ahora bien, $H^i(X,\mathcal O_X)=0$ salvo si $i=0$ o
$i=n=\dim X$, mientras que
$H^0(X,\mathcal O_X)\simeq\mathbb F_q$ (XI~1.5(iii)); por tanto,
\begin{equation*}
\operatorname{card}\bigl(X(\mathbb F_q)\bigr)
\equiv1+(-1)^n
\operatorname{tr}\bigl(\mathfrak F_q\mid H^n(X,\mathcal O_X)\bigr)
\pmod p.
\end{equation*}
Si $X$ es especial, $\mathfrak F_q$ actúa nilpotentemente sobre
$H^n(X,\mathcal O_X)$ y su traza es nula. Esto prueba (1.1.5.1) y la
proposición.

\subsubsection*{Proposición 1.2.}
Para toda potencia $q$ de $p$, si
\begin{equation*}
\sum_{i=1}^{r}d_i\geq n+r+1,
\end{equation*}
existe una intersección completa en $\mathbb P^{n+r}_{\mathbb F_q}$ de
multigrado $(d_1,\dots,d_r)$ que no tiene ningún punto racional sobre
$\mathbb F_q$.

\noindent\underline{Demostración.}
Comencemos por $r=1$. Para todo $d\geq n+2$ debemos hallar una forma
homogénea de grado $d$,

\pageoriginale

\begin{equation}\tag{1.2.1}
H(X_1,\dots,X_{n+2})\in
\mathbb F_q[X_1,\dots,X_{n+2}],
\end{equation}
tal que
\begin{equation}\tag{1.2.2}
\begin{aligned}
&(\alpha_1,\dots,\alpha_{n+2})\in(\mathbb F_q)^{n+2},\qquad
H(\alpha_1,\dots,\alpha_{n+2})=0\quad\Longrightarrow\\
&\hspace{7em}\alpha_1=\cdots=\alpha_{n+2}=0.
\end{aligned}
\end{equation}
Pongamos $\mathbb E=\mathbb F_{q^d}$ y elijamos elementos
$\beta_1,\dots,\beta_{n+2}$ de $\mathbb E$ linealmente independientes
sobre $\mathbb F_q$. Tomemos
\begin{equation}\tag{1.2.3}
H(X_1,\dots,X_{n+2})=
N_{\mathbb E/\mathbb F_q}
\left(\sum_{i=1}^{n+2}\beta_iX_i\right).
\end{equation}

En el caso general, como
\begin{equation}\tag{1.2.4}
\sum_{i=1}^{r}d_i\geq n+r+1,
\end{equation}
podemos poner $a_0=0$ y elegir $r$ enteros
\begin{equation}\tag{1.2.5}
1\leq a_1<a_2<\cdots<a_r=n+r+1
\end{equation}
tales que
\begin{equation}\tag{1.2.6}
d_i\geq a_i-a_{i-1}\qquad(1\leq i\leq r).
\end{equation}

\pageoriginale

Para cada $i=1,\dots,r$, elijamos una forma homogénea de grado $d_i$,
\begin{equation*}
H_i(X_{a_{i-1}+1},\dots,X_{a_i})
\in\mathbb F_q[X_{a_{i-1}+1},\dots,X_{a_i}],
\end{equation*}
en las $a_i-a_{i-1}$ variables indicadas y con la propiedad (1.2.2). La
intersección completa buscada está definida por las formas homogéneas
\begin{equation*}
H_1(X_1,\dots,X_{a_1}),\quad\dots,\quad
H_r(X_{a_{r-1}+1},\dots,X_{a_r}).
\end{equation*}

\subsubsection*{Teorema 1.3.}
Sea $X$ una intersección completa irreducible de dimensión $n\geq1$ sobre
un cuerpo $k$ de exponente característico $p$. Para todo entero
$d\geq n+1$, una sección suficientemente general de $X$ por una
hipersuperficie de grado $d$ no es especial en el sentido de 1.0.

\noindent\underline{Demostración.}
Las secciones $X\mathbin{\cdot}H_u$ de $X$ por hipersuperficies de grado
$d$ están parametrizadas por las rectas de $H^0(X,\mathcal O_X(d))$, es
decir, por
\begin{equation}\tag{1.3.1}
U=\mathbb P\bigl(H^0(X,\mathcal O_X(d))\bigr).
\end{equation}
Sobre $U$ tenemos la «sección universal de $X$ por una hipersuperficie de
grado $d$»
\begin{equation}\tag{1.3.2}
f:X\mathbin{\cdot}H_u\longrightarrow U.
\end{equation}
Sea $U'$ el abierto de los $u\in U$ para los que $\dim(X_u)=n-1$. Por
XI~1.5(i) y 2.5, la restricción
\begin{equation}\tag{1.3.3}
\left.R^{n-1}f_*\mathcal O_{X\cdot H_u}\right|_{U'}
\end{equation}

\pageoriginale

es localmente libre, su formación conmuta con todo cambio de base y su
rango $M$ es al menos uno. Si $p=1$, esto demuestra el teorema para toda
sección por una hipersuperficie de grado al menos $n+1$.

Supongamos, pues, $p>0$. La operación iterada de Hasse--Witt
$\mathfrak F_{p^M}$ se representa de nuevo localmente por una matriz
$M\times M$ con coeficientes en $\mathcal O_U$. Por consiguiente,
\begin{equation}\tag{1.3.4}
\{u\in U'\mid X\mathbin{\cdot}H_u\text{ no es especial}\}
\quad\text{es abierto en }U'.
\end{equation}
Solo queda probar que el conjunto de (1.3.4) no es vacío.

Consideremos primero el caso
\begin{equation}\tag{1.3.5}
k=\mathbb F_q.
\end{equation}
Tras reemplazar $\mathbb F_q$ por una extensión finita, podemos suponer
que
\begin{equation}\tag{1.3.6}
\operatorname{card}\bigl(X(\mathbb F_q)\bigr)>0
\end{equation}
y que
\begin{equation}\tag{1.3.7}
\begin{aligned}
&\text{hay coordenadas homogéneas }X_1,\dots,X_{n+r+1}
 \text{ en }\mathbb P_k^{n+r}\\
&\text{tales que $X$ no corta la subvariedad lineal }
X_1=\cdots=X_{n+1}=0.
\end{aligned}
\end{equation}
Elijamos una forma homogénea de grado $d$,
\begin{equation}\tag{1.3.8}
H(X_1,\dots,X_{n+1})\in
\mathbb F_q[X_1,\dots,X_{n+1}],
\end{equation}
con la propiedad (1.2.2), lo que es posible porque $d\geq n+1$. Sea $H$
también la hipersuperficie de $\mathbb P^{n+r}$ definida por
$H(X_1,\dots,X_{n+1})=0$.

\pageoriginale

Por construcción, $X\mathbin{\cdot}H$ no tiene puntos racionales sobre
$\mathbb F_q$. En vista de (1.3.6), esto muestra además que
$X\not\subset H$; como $X$ es irreducible,
$X\mathbin{\cdot}H$ tiene dimensión $n-1$, luego $H\in U'$. La
Proposición~1.1.5 muestra que $X\mathbin{\cdot}H$ no es especial.

Pasemos al caso general. Examinando las ecuaciones que definen $X$,
podemos suponer que $k$ es de tipo finito sobre $\mathbb F_p$. Existe
entonces un subanillo $B\subset k$, de tipo finito como álgebra sobre
$\mathbb F_p$ y de cuerpo de fracciones $k$, tal que
\begin{equation}\tag{1.3.9}
X\subset\mathbb P_k^{n+r}
\end{equation}
es la fibra genérica de una intersección completa relativa
$\widetilde X\subset\mathbb P_B^{n+r}$ que, fibra por fibra, tiene el
mismo multigrado que $X$ y es geométricamente irreducible.

Escribamos
\begin{equation}\tag{1.3.10}
\rho:\widetilde X\longrightarrow S=\operatorname{Spec}(B)
\end{equation}
para la proyección,
\begin{equation}\tag{1.3.11}
E=\rho_*\mathcal O_{\widetilde X}(d)
\end{equation}
para el $\mathcal O_S$-módulo localmente libre, y
\begin{equation}\tag{1.3.12}
\pi:T=\check{\mathbb P}(E)\longrightarrow S
\end{equation}
para el fibrado de rectas asociado. La fibra genérica de (1.3.12) es
precisamente
\[
U=\mathbb P\bigl(H^0(X,\mathcal O_X(d))\bigr).
\]
Sobre $T$ tenemos la sección universal de
$\widetilde X\subset\mathbb P_B^{n+r}$ por hipersuperficies de grado $d$,

\pageoriginale

denotada
\begin{equation}\tag{1.3.13}
\alpha:\widetilde X\mathbin{\cdot}H_T\longrightarrow T.
\end{equation}
Estas construcciones encajan en el diagrama siguiente de cuadrados
cartesianos:
\begin{equation}\tag{1.3.14}
\begin{tikzcd}[column sep=huge,row sep=large]
\widetilde X\mathbin{\cdot}H_T \arrow[d,"\alpha"]
& X\mathbin{\cdot}H_u=
 \sisi{(\widetilde X\mathbin{\cdot}H_T)\times_U T}
       {(\widetilde X\mathbin{\cdot}H_T)\times_T U}
 \arrow[l]\arrow[d] \\
T \arrow[d,"\pi"]
& U=T\times_S\operatorname{Spec}(k) \arrow[l]\arrow[d] \\
S & \operatorname{Spec}(k)\arrow[l]
\end{tikzcd}
\end{equation}
El $\mathcal O_T$-módulo
\begin{equation}\tag{1.3.15}
R^{n-1}\alpha_*\mathcal O_{\widetilde X\cdot H_T}
\end{equation}
es localmente libre sobre el abierto $T'\subset T$ formado por los $t$
para los que $\dim\alpha^{-1}(t)=n-1$, y su formación conmuta allí con
todo cambio de base. Como $T$ y $U$ tienen el mismo punto genérico, para
mostrar que (1.3.4) no es vacío basta probar que el abierto de $T'$
\begin{equation}\tag{1.3.16}
\bigl\{t\in T'\mid\widetilde X_{\pi(t)}\mathbin{\cdot}H_t
\text{ no es especial}\bigr\}
\end{equation}
no es vacío. Pero el caso (1.3.5), ya demostrado, muestra que, para todo
punto cerrado $s\in S$, la fibra $\pi^{-1}(s)$ corta al abierto (1.3.16).

\pageoriginale

Esto demuestra el teorema.

\subsubsection*{Teorema 1.4.}
Sea $X\subset\mathbb P^N$ un esquema proyectivo Cohen--Macaulay,
equidimensional de dimensión $n\geq1$, sobre un cuerpo $k$. Existe un
entero $d_0$ tal que, para todo $d\geq d_0$, si
$i:X\mathbin{\cdot}H\longrightarrow X$ denota la inclusión de una sección
por una hipersuperficie general de grado $d$, entonces la acción de
\begin{equation}\tag{1.4.0}
\mathfrak F_p\quad\text{sobre}\quad
\frac{H^{n-1}(X\mathbin{\cdot}H,\mathcal O_{X\cdot H})}
{i^*H^{n-1}(X,\mathcal O_X)}
\end{equation}
no es nilpotente.

\noindent\underline{Demostración.}
Como antes, nos reducimos al caso $k=\mathbb F_q$ y a hallar un solo $H$
---para cada grado fijo suficientemente grande para el que
$\dim H^{n-1}(X\mathbin{\cdot}H,\mathcal O_{X\cdot H})$ sea independiente
de $H$; véase (1.4.4)--- que verifique (1.4.0). Ampliando $q$, podemos
suponer que las componentes irreducibles $X_i$ de $X$ son
geométricamente irreducibles, que
\begin{equation}\tag{1.4.1}
\operatorname{card}\bigl(X_i(\mathbb F_q)\bigr)\geq2,
\end{equation}
y que
\begin{equation}\tag{1.4.2}
\begin{aligned}
&\text{hay coordenadas proyectivas }X_0,\dots,X_N\text{ tales que}\\
&\qquad X\text{ no corta la subvariedad lineal }
X_0=\cdots=X_n=0.
\end{aligned}
\end{equation}

\pageoriginale

Por un teorema de anulación de Serre ([1], o SGA~2 XII~1.4), existe un
entero $d_1$ tal que, para todo $d\geq d_1$,
\begin{equation}\tag{1.4.3}
H^i(X,\mathcal O_X(-d))=0\qquad\text{si }i\ne n=\dim X.
\end{equation}
Así, si $H$ es una hipersuperficie de grado $d\geq d_1$, (1.4.3) da
\begin{equation}\tag{1.4.4}
\begin{aligned}
H^i(X,\mathcal O_X)&\xrightarrow{\ \sim\ }
H^i(X\mathbin{\cdot}H,\mathcal O_{X\cdot H})
&&\text{si }i\leq n-2,\\[1ex]
H^{n-1}(X,\mathcal O_X)&\lhook\joinrel\longrightarrow
H^{n-1}(X\mathbin{\cdot}H,\mathcal O_{X\cdot H})
&&\text{inyectivamente}.
\end{aligned}
\end{equation}

Demostremos el teorema con $d_0=\max(N+1,d_1)$. Para cada $d\geq d_0$
debemos hallar una hipersuperficie $H$ de grado $d$, definida sobre
$\mathbb F_q$, que no contenga ningún $X_i$ y para la que
\begin{equation}\tag{1.4.5}
\operatorname{tr}\!\left(\mathfrak F_q\ \middle|\
\frac{H^{n-1}(X\cdot H,\mathcal O_{X\cdot H})}
{H^{n-1}(X,\mathcal O_X)}\right)\ne0.
\end{equation}
Por la fórmula de las trazas (XXII~3.2.1), módulo $p$ tenemos
\begin{equation}\tag{1.4.6}
\begin{aligned}
\operatorname{card}\bigl((X\cdot H)(\mathbb F_q)\bigr)
&\equiv\sum_{i=0}^{n-1}(-1)^i
\operatorname{tr}\bigl(\mathfrak F_q\mid
H^i(X\cdot H,\mathcal O_{X\cdot H})\bigr)\\
&\stackrel{(1.4.4)}{\equiv}
\sum_{i=0}^{n-1}(-1)^i
\operatorname{tr}\bigl(\mathfrak F_q\mid H^i(X,\mathcal O_X)\bigr)\\
&\quad+(-1)^{n-1}\operatorname{tr}\!\left(
\mathfrak F_q\ \middle|\
\frac{H^{n-1}(X\cdot H,\mathcal O_{X\cdot H})}
{H^{n-1}(X,\mathcal O_X)}\right).
\end{aligned}
\end{equation}

\pageoriginale

Sea $a\in\mathbb Z/p\mathbb Z$ la cantidad
\[
a=\sum_{i=0}^{n-1}(-1)^i
\operatorname{tr}\bigl(\mathfrak F_q\mid H^i(X,\mathcal O_X)\bigr).
\]
Basta, pues, para todo $d\geq d_0$, hallar una hipersuperficie $H$ de
grado $d$, definida sobre $\mathbb F_q$, que no contenga $X$ y tal que
\begin{equation}\tag{1.4.7}
\operatorname{card}\bigl((X\cdot H)(\mathbb F_q)\bigr)
\not\equiv a\pmod p.
\end{equation}

Hay dos casos. Si $a\ne0$, basta tomar un $H$ sobre $\mathbb F_q$ para el
que $(X\cdot H)(\mathbb F_q)$ sea vacío, como en (1.3.8). Tal $H$ verifica
(1.4.7) y no contiene ningún $X_i$, pues
$\operatorname{card}(X_i(\mathbb F_q))\geq2$.

Si $a=0$, basta hallar un $H$ tal que
$\operatorname{card}((X\cdot H)(\mathbb F_q))=1$. Entonces (1.4.7) se
verifica y $X_i\not\subset H$. Elijamos coordenadas proyectivas
$X_0,\dots,X_N$ tales que $(1,0,\dots,0)$ sea un punto de $X$. Tomemos
$H$ de ecuación $H(X_1,\dots,X_N)=0$, donde $H$ es una forma homogénea
que verifica
\begin{equation}\tag{1.4.8}
\alpha_1,\dots,\alpha_N\in\mathbb F_q,\quad
H(\alpha_1,\dots,\alpha_N)=0
\quad\Longrightarrow\quad
\alpha_1=\cdots=\alpha_N=0.
\end{equation}
La hipersuperficie definida por $H$ en $\mathbb P^N$ tiene entonces como
único punto sobre $\mathbb F_q$ a $(1,0,\dots,0)$; a fortiori,
$X\cdot H$ solo tiene ese punto. Esto demuestra el teorema.

\pageoriginale

\section*{2. El conivel sobre un cuerpo finito de definición}

\subsection*{2.0.}
Sea $X$ un esquema separado de tipo finito sobre $\mathbb F_q$. Para todo
primo $\ell\ne\operatorname{char}(\mathbb F_q)$, escribamos
\begin{equation}\tag{2.0.1}
H_c^i(\overline X,\mathbb Q_\ell)
\end{equation}
para los grupos de cohomología con soporte propio de
$\overline X=X\otimes_{\mathbb F_q}\overline{\mathbb F}_q$. Son espacios
vectoriales de dimensión finita sobre $\mathbb Q_\ell$ en los que actúa
$\operatorname{Gal}(\overline{\mathbb F}_q/\mathbb F_q)$. Sea
\begin{equation}\tag{2.0.2}
\varphi_q\in\operatorname{Gal}(\overline{\mathbb F}_q/\mathbb F_q)
\end{equation}
el automorfismo de Frobenius. Se sabe (SGA~5 XIII~3.2.1) que la función
zeta de $X$ (XXII~3.0.1) viene dada por
\begin{equation}\tag{2.0.3}
Z(t,X/\mathbb F_q)=
\prod_{i\geq0}
\det\!\left(1-t\varphi_q^{-1}\mid
H_c^i(\overline X,\mathbb Q_\ell)\right)^{(-1)^{i+1}}.
\end{equation}
Deligne ha demostrado recientemente, generalizando un resultado de
Washnitzer para $X$ proyectivo liso, el teorema siguiente, que se demuestra
en el apéndice (5.5.3).

\subsubsection*{Teorema 2.1.}
Los valores propios de $\varphi_q^{-1}$ sobre
$H_c^i(\overline X,\mathbb Q_\ell)$ son enteros algebraicos; si
$i>\dim X$, son divisibles, como enteros algebraicos, por
$q^{i-\dim X}$.

\pageoriginale

En particular se obtiene, tautológicamente:

\medskip
\noindent\textbf{(2.1 bis).}
Los valores propios de $\varphi_q^{-1}$ sobre
$H_c^i(\overline X,\mathbb Q_\ell(-j))$ son enteros algebraicos divisibles
por $q^j$.

\subsubsection*{Corolario 2.2.}
Sea $X$ propio, liso y geométricamente conexo sobre $\mathbb F_q$.
Entonces los valores propios de $\varphi_q^{-1}$ sobre
$N^jH^q(\overline X,\mathbb Q_\ell)$ (notación de XX~2.0.6) son enteros
algebraicos divisibles por $q^j$.

\noindent\underline{Demostración.}
La afirmación es invariante por una extensión finita
$\mathbb F_{q^M}/\mathbb F_q$, que equivale a reemplazar los valores
propios considerados por sus potencias $M$-ésimas. Tras tal cambio de
base, la definición proporciona un número finito de esquemas $T_i$
proyectivos, lisos y geométricamente conexos sobre $\mathbb F_q$, de
dimensiones $d_i$, y ciclos algebraicos $Z_i$ sobre
$T_i\times_{\mathbb F_q}X$ de codimensión $j+d_i$, tales que el
homomorfismo de $\operatorname{Gal}(\overline{\mathbb F}_q/\mathbb F_q)$-
módulos
\begin{equation*}
\bigoplus_i H^{q-2j}(\overline T_i,\mathbb Q_\ell(-j))
\xrightarrow{\ \oplus_i\operatorname{Cl}(Z_i)\ }
N^jH^q(\overline X,\mathbb Q_\ell)
\end{equation*}
es sobreyectivo. La conclusión se deduce de (2.1 bis).

\pageoriginale

\section*{3. Aplicación a las intersecciones completas}

\subsection*{3.0.}
Sea $X$ una intersección completa lisa de dimensión $n\geq1$ en
$\mathbb P^N$ sobre un cuerpo $k$. Para
$\ell\ne\operatorname{char}(k)$, Lefschetz débil y la dualidad dan
(XI~1.6)
\begin{equation}\tag{3.0.1}
H^i(\overline X,\mathbb Q_\ell)=
\begin{cases}
0 & \text{si $i$ es impar e $i\ne n$},\\[1.5ex]
\mathbb Q_\ell\!\left(-\dfrac{i}{2}\right)
  & \text{si $i$ es par, $i\leq2n$ e $i\ne n$},\\[1.5ex]
H^n(\overline X,\mathbb Q_\ell) & \text{si $i=n$}.
\end{cases}
\end{equation}
Para $k=\mathbb F_q$ se obtiene la fórmula siguiente para la función zeta:
\begin{equation}\tag{3.0.2}
Z(t,X/\mathbb F_q)\prod_{i=0}^{n}(1-q^it)
=\det\!\left(1-t\varphi_q^{-1}\mid
\mathrm{Prim}^n(\overline X,\mathbb Q_\ell)\right)^{(-1)^{n+1}}.
\end{equation}
Aquí se ha utilizado la descomposición primitiva y (3.0.1).

\subsubsection*{Corolario 3.0.3.}
Bajo las hipótesis de 3.0, con $k=\mathbb F_q$, se tiene la congruencia
módulo $p$
\[
\det\!\left(1-t\varphi_q^{-1}\mid
H^n(\overline X,\mathbb Q_\ell)\right)
\equiv
\det\!\left(1-t\mathfrak F_q\mid H^n(X,\mathcal O_X)\right).
\]

\noindent\underline{Demostración.}
Por la fórmula de congruencia (XXII~3.2.1), y puesto que
\begin{equation}\tag{3.0.4}
H^i(X,\mathcal O_X)\simeq
\begin{cases}
\mathbb F_q & i=0,\\[1.5ex]
0 & 1\leq i\leq n-1,\\[1.5ex]
0 & i\geq n+1,
\end{cases}
\end{equation}

\pageoriginale

tenemos
\begin{equation}\tag{3.0.5}
Z(t,X/\mathbb F_q)\equiv
(1-t)^{-1}
\det\!\left(1-t\mathfrak F_q\mid H^n(X,\mathcal O_X)\right)^{(-1)^{n+1}}
\pmod p.
\end{equation}
Por consiguiente, (3.0.5) y (3.0.2) dan
\begin{equation}\tag{3.0.6}
\det\!\left(1-t\varphi_q^{-1}\mid
\mathrm{Prim}^n(\overline X,\mathbb Q_\ell)\right)
\equiv
\det\!\left(1-t\mathfrak F_q\mid H^n(X,\mathcal O_X)\right)
\pmod p.
\end{equation}
La conclusión buscada resulta de
\begin{equation}\tag{3.0.7}
H^n(\overline X,\mathbb Q_\ell)\simeq
\begin{cases}
\mathrm{Prim}^n(\overline X,\mathbb Q_\ell) & \text{si $n$ es impar},\\[2mm]
\mathrm{Prim}^n(\overline X,\mathbb Q_\ell)
\oplus\mathbb Q_\ell\!\left(-\dfrac n2\right) & \text{si $n$ es par}.
\end{cases}
\end{equation}

\subsubsection*{Corolario 3.1.}
Sea $X$ una intersección completa lisa de dimensión $n\geq1$ sobre
$\mathbb F_q$. Si
\begin{equation}\tag{3.1.0}
H^n(\overline X,\mathbb Q_\ell)=N^1H^n(\overline X,\mathbb Q_\ell),
\end{equation}
entonces
\begin{equation}\tag{3.1.1}
\mathfrak F_p\quad\text{es nilpotente sobre}\quad H^n(X,\mathcal O_X).
\end{equation}

\noindent\underline{Demostración.}
Por el Corolario~2.2, (3.1.0) implica
\begin{equation}\tag{3.1.2}
\det\!\left(1-t\varphi_q^{-1}\mid
H^n(\overline X,\mathbb Q_\ell)\right)
\equiv1\pmod p
\qquad\text{(incluso módulo $q$)}.
\end{equation}
Junto con el Corolario~3.0.3, esto da
\begin{equation}\tag{3.1.3}
\det\!\left(1-t\mathfrak F_q\mid H^n(X,\mathcal O_X)\right)=1.
\end{equation}

\pageoriginale

Por tanto, $\mathfrak F_q$ es nilpotente sobre $H^n(X,\mathcal O_X)$, y
también lo es $\mathfrak F_p$. Esto demuestra el corolario.

\subsubsection*{Proposición 3.2.}
Sea $X$ una intersección completa lisa de dimensión $n\geq1$ sobre un
cuerpo $k$. Si
\begin{equation}\tag{3.2.1}
H^n(\overline X,\mathbb Q_\ell)=N^1H^n(\overline X,\mathbb Q_\ell),
\end{equation}
entonces $X$ es especial, es decir,
\begin{equation}\tag{3.2.2}
\mathfrak F_p\quad\text{es nilpotente sobre}\quad H^n(X,\mathcal O_X).
\end{equation}

\noindent\underline{Demostración.}
Podemos suponer que $k$ es de tipo finito sobre $\mathbb F_p$ (el caso
$p=0$ es trivial), y que existe una $\mathbb F_p$-subálgebra $B$ de $k$,
de tipo finito, tal que $X/k$ sea la fibra genérica de una intersección
completa relativa lisa $\widetilde X$:
\begin{equation}\tag{3.2.3}
\begin{tikzcd}[column sep=4.5em,row sep=3em]
\widetilde X \arrow[d,"\pi"'] &
X=\sisi{\widetilde X/_{\eta}}{\widetilde X_\eta}
 \arrow[l]\arrow[d] \\
\operatorname{Spec}(B)=S & \eta=\operatorname{Spec}(k)\arrow[l].
\end{tikzcd}
\end{equation}
El $B$-módulo $R^n\pi_*\mathcal O_{\widetilde X}$ es localmente libre y
su formación conmuta con todo cambio de base. Por consiguiente, el conjunto
\begin{equation}\tag{3.2.4}
\{s\in S\mid X_s\text{ es especial}\}
\end{equation}

\pageoriginale

es cerrado en $S$.

Por XX~2.3, tras encoger $S$ podemos suponer que
\begin{equation}\tag{3.2.5}
\forall s\in S,\qquad
H^n(X_{\bar s},\mathbb Q_\ell)=N^1H^n(X_{\bar s},\mathbb Q_\ell).
\end{equation}
Por el Corolario~3.1, $X_s$ es especial para todo punto cerrado $s$ de
$S$. Se sigue que el cerrado (3.2.4) es todo $S$ y, en particular, que
$X=\widetilde X_\eta$ es especial.

\section*{4. Conclusiones}

Reuniendo el Teorema~1.1 y la Proposición~3.2 se obtiene:

\subsubsection*{Teorema 4.1.}
Fijemos una característica $p$, una dimensión $n$ y un multigrado
$(d_1,\ldots,d_r)$. Sea
\begin{equation}\tag{4.1.0}
X_V\longrightarrow V
\end{equation}
la intersección completa lisa universal de dimensión $n$ y multigrado
$(d_1,\ldots,d_r)$ en característica $p$. Si
\begin{equation}\tag{4.1.1}
\sum_{i=1}^{r}d_i\geq n+r+1,
\end{equation}
entonces existe un abierto denso $U\subset V$ tal que
\begin{equation}\tag{4.1.2}
\text{para todo }u\in U,\qquad
H^n(X_{\bar u},\mathbb Q_\ell)\ne N^1H^n(X_{\bar u},\mathbb Q_\ell).
\end{equation}

Reuniendo de igual modo el Teorema~1.3 y la Proposición~3.2 se obtiene:

\pageoriginale

\subsubsection*{Teorema 4.2.}
Sea $X$ una intersección completa lisa de dimensión $n\geq1$ sobre un
cuerpo $k$. Fijemos un entero
\begin{equation}\tag{4.2.1}
d\geq n+1.
\end{equation}
Sea
\begin{equation}\tag{4.2.2}
X\mathbin{\cdot}H_u\longrightarrow U
\end{equation}
la familia universal de secciones de $X$ por hipersuperficies de grado
$d$ que son lisas de dimensión $n-1$. Existe un abierto denso $T\subset U$
tal que
\begin{equation}\tag{4.2.3}
\text{para todo }t\in T,\qquad
H^{n-1}(X\mathbin{\cdot}H_{\bar t},\mathbb Q_\ell)
\ne N^1H^{n-1}(X\mathbin{\cdot}H_{\bar t},\mathbb Q_\ell).
\end{equation}

\section*{Bibliografía}
\begin{enumerate}[label={\textup{[\arabic*]}},leftmargin=2.6em]
\item J.-P. Serre, \emph{Faisceaux algébriques cohérents},
  \emph{Annals of Mathematics} \textbf{61} (1955), 197--278.
\item J.-P. Serre, \emph{Sur la topologie des variétés algébriques en
  caractéristique $p$}, en \emph{Symposium Internacional de Topología
  Algebraica}, México, 1958.
\end{enumerate}

\pageoriginale

\section*{5. Apéndice: teoremas de integralidad}
\begin{center}
por P. Deligne
\end{center}

\subsection*{5.0.}
Sea $p$ un número primo, $q$ una potencia de $p$, $\mathbb F_q$ un cuerpo
de $q$ elementos, $E$ el completado de un cuerpo de números en un lugar
que divide un primo $\ell\ne p$, y $T$ un conjunto de números primos (por
ejemplo, $T=\varnothing$).

Un elemento $\alpha$ de una clausura algebraica de $E$ se llama
\emph{$T$-entero} si es algebraico sobre $\mathbb Q$ y entero sobre
$\mathbb Z[(1/r)_{r\in T}]$. Un endomorfismo $F$ de un espacio vectorial
de dimensión finita sobre $E$ se llama $T$-entero si sus valores propios
son $T$-enteros.

Sea $k$ un cuerpo finito de $q'$ elementos, $\bar k$ una clausura
algebraica de $k$ y
\[
\mathfrak G=\operatorname{Gal}(\bar k/k).
\]
Denotemos por $f\in\mathfrak G$ la sustitución de Frobenius
\[
f(x)=x^{q'},
\]
y por $F$ el Frobenius «geométrico», inverso de $f$ en $\mathfrak G$.
Una representación $\rho$ de $\mathfrak G$ sobre un espacio vectorial de
dimensión finita sobre $E$ se llama $T$-entera si $\rho(F)$ es $T$-entero.

Si $X$ es un esquema de tipo finito sobre $\mathbb F_q$ y $x$ es un punto
cerrado de $X$, denotamos por $F_x$ el elemento de Frobenius geométrico del
grupo de Galois de $k(x)$. Si $\bar x$ es un punto geométrico de $X$ sobre
$x$ y $\underline G$ es un $E$-haz sobre $X$, entonces $F_x$ actúa, por
transporte de estructura, sobre la fibra $G_{\bar x}$.

\subsubsection*{Definición 5.1.}
Un $E$-haz constructible $\underline G$ sobre un esquema $X$ de tipo
finito sobre $\mathbb F_q$ se llama $T$-entero si, para todo punto
geométrico $\bar x$ de $X$ sobre un punto cerrado $x$, $F_x$ es un
automorfismo $T$-entero de $G_{\bar x}$.

\pageoriginale

Si $a:X\longrightarrow\operatorname{Spec}(\mathbb F_q)$ es un esquema
de tipo finito sobre $\mathbb F_q$ y $\underline G$ un $E$-haz
constructible sobre $X$, escribimos $\underline H^*(X,\underline G)$ para
los $E$-haces $R^*a_*\underline G$ sobre
$\operatorname{Spec}(\mathbb F_q)_{\acute e t}$. Si
$\overline{\mathbb F}_q$ es una clausura algebraica de $\mathbb F_q$ y
$\overline X=X\times_{\mathbb F_q}\overline{\mathbb F}_q$, estos haces se
identifican con los grupos $H^*(\overline X,\underline G)$ considerados
como módulos galoisianos, es decir, provistos de una acción del Frobenius
geométrico $F$. Adoptamos la misma convención para la cohomología con
soporte propio.

\subsubsection*{Lema 5.2.1.}
Sea $\underline G$ un $E$-haz constructible sobre un esquema $X$ de tipo
finito sobre $\mathbb F_q$, de dimensión a lo sumo $n$.
\begin{enumerate}[label=(\roman*)]
\item Existe un cerrado $Y$ de $X$, de dimensión $0$, tal que la
  aplicación evidente
  \[
  \underline H^0(X,\underline G)\longrightarrow
  \underline H^0(Y,\underline G)
  \]
  es inyectiva.
\item Si $X$ es separado y $U$ es un abierto de $X$ cuyo complemento $Z$
  tiene dimensión menor que $n$, existe un cerrado $Y$ de $U$, de
  dimensión $0$, y una aplicación sobreyectiva
  \[
  \underline H^0(Y,\underline G)(-n)\longrightarrow
  \underline H_c^{2n}(X,\underline G).
  \]
\end{enumerate}

\noindent\underline{Demostración.}
La afirmación (i) es evidente. Demostremos (ii). Tras reemplazar $X$ por
$X_{\mathrm{red}}$ y encoger $U$, podemos suponer que $U$ es liso y que
$\underline G|_U$ es constante torcido. Como $\dim Z<n$, tenemos
\[
0=\underline H_c^{2n-1}(Z,\underline G)
\longrightarrow\underline H_c^{2n}(U,\underline G)
\xrightarrow{\sim}\underline H_c^{2n}(X,\underline G)
\longrightarrow\underline H_c^{2n}(Z,\underline G)=0,
\]
y la dualidad de Poincaré (SGA~4, Exposé~XVIII, y [2]) da
\[
\underline H_c^{2n}(U,\underline G)\simeq
\underline H^0(U,\underline G^\vee(n))^\vee.
\]
La conclusión se deduce de (i).

\pageoriginale

\subsubsection*{Teorema 5.2.2.}
Sea $X$ un esquema separado de tipo finito sobre $\mathbb F_q$, de
dimensión a lo sumo $n$, y sea $\underline G$ un $E$-haz constructible
$T$-entero. Entonces los endomorfismos $F$ y $q^{n-i}F$ de
$\underline H_c^i(X,\underline G)$ son $T$-enteros.

\noindent\underline{Demostración.}
El caso $\dim X\leq0$ es trivial, y los casos $i=0$ e $i=2n$ resultan del
Lema~5.2.1. Sea
\[
Z(X,\underline G,t)=
\prod_{x\in X}\det\!\left(1-F_xt^{\deg(x)}\,;\,G_x\right)^{-1}
\in E[[t]].
\]
Por SGA~5, Exposé~XV, 3.2, se tiene
\begin{equation}\tag{5.2.3}
Z(X,\underline G,t)=
\prod_i\det\!\left(1-Ft\,;\,
\underline H_c^i(X,\underline G)\right)^{(-1)^{i+1}}.
\end{equation}

Demostremos primero el teorema para $n=1$. Entonces
$\underline H_c^i(X,\underline G)=0$ para $i>2$ (SGA~4, Exposé~X, 4.3,
o Exposé~XVII, 5.2.8.1). Por hipótesis, los coeficientes de
$Z(X,\underline G,t)$ son $T$-enteros. Por los casos $i=0$ e $i=2n=2$
ya demostrados, el producto de $Z(X,\underline G,t)$ por
\[
\det\!\left(1-Ft\,;\,\underline H_c^0(X,\underline G)\right)
\det\!\left(1-Ft\,;\,\underline H_c^2(X,\underline G)\right)
\]
tiene aún coeficientes $T$-enteros. Este producto es precisamente
$\det(1-Ft\,;\,\underline H_c^1(X,\underline G))$, lo que demuestra el
teorema para $n=1$.

Procedamos por inducción sobre $n=\dim X$. Sea $U$ un abierto de Zariski
de $X$ cuyo complemento $Z$ tiene dimensión menor que $n$, y supongamos
que existe un morfismo $f:U\longrightarrow Y$ hacia una curva, de fibras
de dimensión a lo sumo $n-1$:

\pageoriginale

\begin{center}
\begin{tikzcd}
U \arrow[r,hook,"j"] \arrow[d,"f"'] & X & Z \arrow[l,hook',"i"'] \\
Y
\end{tikzcd}
\qquad .
\end{center}
La sucesión exacta
\[
0\longrightarrow j_!j^*\underline G\longrightarrow\underline G
\longrightarrow i_!i^*\underline G\longrightarrow0
\]
induce sucesiones exactas
\[
\underline H_c^i(U,\underline G)\longrightarrow
\underline H_c^i(X,\underline G)\longrightarrow
\underline H_c^i(Z,\underline G).
\]
Por inducción, basta comprobar el teorema para $\underline G$ sobre $U$.
Aplicando la hipótesis de inducción a las fibras de $f$, junto con el
cambio de base propio, se ve que los haces $R^if_!\underline G$ y
$R^if_!\underline G(i-(n-1))$ sobre $Y$ son $T$-enteros. La sucesión
espectral de Leray para $f$ es
\[
E_2^{pq}=\underline H_c^p(Y,R^qf_!\underline G)
\Longrightarrow\underline H_c^{p+q}(U,\underline G).
\]
Por lo anterior y por el teorema ya demostrado para $Y$, tanto $E_2^{pq}$
como $E_2^{pq}(p+q-n)$ son $T$-enteros. Se sigue el teorema.

\subsubsection*{Corolario 5.3.}
Sea $X$ un esquema de tipo finito sobre $\mathbb F_q$, de dimensión a lo
sumo $1$, sea $U$ un abierto de $X$, sea $j:U\hookrightarrow X$ la
inclusión, y sea $\underline G$ un $E$-haz constructible sobre $U$.
\begin{enumerate}[label=(\alph*)]
\item Si $\underline G$ es $T$-entero, entonces $j_*\underline G$ es
  $T$-entero.
\item Si $G_x^\vee$ es $T$-entero para todo $x\in U$, entonces
  $(j_*\underline G)_x^\vee$ es $T$-entero para todo $x\in X$.
\end{enumerate}

\pageoriginale

\noindent\underline{Demostración.}
Nos reducimos primero al caso en que $X$ es afín, y después al caso en que
$X$ es proyectivo. Definamos $\underline H$ por la sucesión exacta
\[
0\longrightarrow j_!\underline G\longrightarrow j_*\underline G
\longrightarrow\underline H\longrightarrow0.
\]
Si $Y$ es el soporte finito de $\underline H$, la sucesión de cohomología
correspondiente da una sucesión exacta
\[
\underline H^0(U,\underline G)\longrightarrow
\underline H^0(Y,\underline H)\longrightarrow
\underline H_c^1(U,\underline G).
\]
El Lema~5.2.1 y el Teorema~5.2.2 muestran que los dos términos exteriores
son $T$-enteros. Por tanto, $\underline H^0(Y,\underline H)$ es $T$-entero,
lo que demuestra (a).

Sea $X'$ la normalización de la parte puramente unidimensional de
$X_{\mathrm{red}}$, y elijamos un abierto $U'$ de $U$, denso en el lugar
liso unidimensional de $U_{\mathrm{red}}$, tal que
$\underline G|_{U'}$ sea constante torcido:
\begin{center}
\begin{tikzcd}
U' \arrow[r,hook,"j'"'] \arrow[d,hook] & X' \arrow[d,hook,"p"] & {} \\
U \arrow[r,hook,"j"'] & X & Y \arrow[l,hook']
\end{tikzcd}
\qquad .
\end{center}
Entonces
\[
(p_*j'_*\underline G)|_Y\simeq(j_*\underline G)|_Y.
\]
Basta tratar el caso en que $X$ es liso de dimensión uno y
$\underline G|_U$ es constante torcido. El haz $\underline G^\vee$ es
entonces $T$-entero, de modo que ya sabemos que
$(j_*\underline G^\vee)_x$ es $T$-entero para todo $x\in X$.

\pageoriginale

Sea $\mathfrak G_x$ el grupo de Galois del cuerpo de fracciones de
$\widehat{\mathcal O}_{X,x}$ y sea $\mathfrak I_x$ el grupo de inercia.
El haz $\underline G$ define una representación de $\mathfrak G_x$,
denotada también por $\underline G$. El
$\mathfrak G_x/\mathfrak I_x=
\operatorname{Gal}(\overline{k(x)}/k(x))$-módulo
$(j_*\underline G^\vee)_x$ (respectivamente,
$(j_*\underline G)_x^\vee$) se identifica con el módulo de invariantes
(respectivamente, de coinvariantes) en $\underline G^\vee$. Queda por
demostrar:

\subsubsection*{Lema 5.3.1.}
Sea $K$ un cuerpo local de cuerpo residual $\mathbb F_q$, sea $G$ su
grupo de Galois, sea $I$ el grupo de inercia y sea $\rho$ una
representación continua de $G$ sobre un espacio vectorial $V$ de
dimensión $n$ sobre $E$ (cf.~5.0). Sean $V^I$ y $V_I$, respectivamente,
el espacio de invariantes y el de coinvariantes bajo $I$, considerados
como $G/I$-módulos. Pongamos $r=\dim_E V^I$, y sean
$\alpha_1,\ldots,\alpha_r$ los valores propios de $F$ sobre $V^I$.
Entonces existen enteros $m_i\geq0$ tales que los valores propios de $F$
sobre $V_I$ son
\[
\alpha_iq^{m_i}\qquad(1\leq i\leq r).
\]

\noindent\underline{Demostración.}
Sea $I'$ el mayor subgrupo de $I$ de orden primo con $\ell$. El grupo
$\rho(I')$ es finito, pues es un subgrupo compacto de
$\operatorname{GL}(V)$ de orden primo con $\ell$. Así,
\[
V_I\simeq\bigl(V^{\rho(I')}\bigr)_I.
\]
Reemplazando $V$ por $V^{I'}$, podemos suponer que $I'$ actúa trivialmente
sobre $V$. Como $G/I$-módulos se tiene
\[
I/I'\simeq\mathbb Z_\ell(1).
\]
Sea $T$ un generador topológico de $I/I'$. Para cada $\alpha$, sea
$V^{(\alpha)}$ el subespacio propio generalizado de $T$ sobre $V$ de
valor propio $\alpha$. Entonces
\[
V=\bigoplus_\alpha V^{(\alpha)},
\]
y la descomposición
\[
V=V^{(1)}\oplus\bigoplus_{\alpha\ne1}V^{(\alpha)}
\]
es estable bajo $G$ y compatible con la dualidad.

\pageoriginale

Podemos, por tanto, reemplazar $V$ por $V^{(1)}$ y suponer que $I$ actúa
unipotentemente.

\subsubsection*{Lema 5.3.2.}
Sea $V$ un espacio vectorial de dimensión finita sobre un cuerpo $k$ de
característica $0$, sea $I$ un grupo aditivo uniparamétrico, sea
$\rho:I\longrightarrow\operatorname{GL}(V)$ una representación y sea $N$
un generador infinitesimal de $I$. Pongamos
\[
V^{I,k}=V^I\cap\operatorname{Im}(N^k)
\]
y
\[
V_{I,k}=\operatorname{Im}\!\left(\operatorname{Ker}(N^{k+1})
\longrightarrow V_I\right).
\]
Entonces $N^k$ induce un isomorfismo
\[
N^k:V_{I,k}/V_{I,k-1}\xrightarrow{\sim}
V^{I,k}/V^{I,k+1},
\]
y define isomorfismos canónicos
\[
V_{I,k}/V_{I,k-1}\otimes I^{\otimes k}
\simeq V^{I,k}/V^{I,k+1}.
\]

\noindent\underline{Demostración.}
Por medio de $N$, las representaciones $V$ de $I$ se identifican con los
$k[T]$-módulos anulados por una potencia de $T$. Basta comprobar el
Lema~5.3.2 cuando $V$ corresponde al módulo indescomponible
$k[T]/T^k$, donde la afirmación es inmediata.

Terminemos la demostración del Lema~5.3.1. El Lema~5.3.2 proporciona
filtraciones creciente y decreciente $V_{I,k}$ y $V^{I,k}$ de $V_I$ y
$V^I$, junto con isomorfismos de $G/I$-módulos
\[
V^{I,k}/V^{I,k+1}\simeq
\bigl(V_{I,k}/V_{I,k-1}\bigr)(k).
\]
Se sigue la afirmación.

\pageoriginale

\subsubsection*{Teorema 5.4.}
Sea $X$ un esquema separado de tipo finito sobre $\mathbb F_q$, de
dimensión $n$, y sea $\underline G$ un $E$-haz constructible sobre $X$.
Supongamos que, para todo $x\in X$, el endomorfismo $F_x$ de $G_x^\vee$
es $T$-entero. Entonces los endomorfismos $q^iF^\vee$ y $q^nF^\vee$ de
$\underline H_c^i(X,\underline G)^\vee$ son $T$-enteros.

\noindent\underline{Demostración.}
Como en el Teorema~5.2.2, nos reducimos al caso en que $X$ es una curva
lisa y $\underline G$ es constante torcido sobre $X$. Sea $j$ la
inclusión de $X$ en una curva proyectiva lisa $X'$ que la compactifica.

Por dualidad de Poincaré, $\underline H^i(X',j_*\underline G)$ es dual de
$\underline H^{2-i}(X',j_*\underline G^\vee(1))$. Por el
Corolario~5.3(a), el Teorema~5.2.2 se aplica a
$j_*\underline G^\vee$:
\begin{align*}
\underline H^0(X',j_*\underline G^\vee)
&\simeq\underline H^2(X',j_*\underline G)^\vee(-1)
&&\text{es $T$-entero},\\
\underline H^1(X',j_*\underline G^\vee)
&\simeq\underline H^1(X',j_*\underline G)^\vee(-1)
&&\text{es $T$-entero},\\
\underline H^2(X',j_*\underline G^\vee)(1)
&\simeq\underline H^0(X',j_*\underline G)^\vee
&&\text{es $T$-entero}.
\end{align*}
Esto demuestra el teorema para $j_*\underline G$.

Definamos $\underline H$ por la sucesión exacta
\[
0\longrightarrow j_!\underline G\longrightarrow j_*\underline G
\longrightarrow\underline H\longrightarrow0.
\]
Si $Y$ es el soporte de $\underline H$, la sucesión exacta larga de
cohomología proporciona
\begin{center}
\begin{tikzcd}[row sep=large,column sep=large]
{} & \underline H_c^0(X,\underline G)
 \arrow[r,hook] & \underline H^0(X',j_*\underline G) & ,\\
\underline H^0(Y,\underline H) \arrow[r] &
 \underline H_c^1(X,\underline G) \arrow[r] &
 \underline H^1(X',j_*\underline G) & ,\\
{} & \underline H_c^2(X,\underline G)
 \arrow[r,"\sim"] & \underline H^2(X',j_*\underline G) & .
\end{tikzcd}
\end{center}

\pageoriginale

Por consiguiente, el teorema para $(X,\underline G)$ resulta del teorema
para $(X',j_*\underline G)$.

\subsection*{5.5.}
El método anterior de fibración en curvas proporciona asimismo las
estimaciones de Lang--Weil [3] para los valores absolutos complejos de los
valores propios de Frobenius.

Sea $X$ un esquema de tipo finito sobre $\mathbb F_q$, sea $\underline G$
un $E$-haz constructible sobre $X$ y sea $\lambda\in\mathbb R$.
Supongamos que, para todo punto $x\in X$, con
$\#k(x)=q_x$, los valores propios de $F_x$ sobre $G_x$ son números
algebraicos $\alpha$ cuyos valores absolutos complejos verifican
\[
|\alpha|\leq q_x^\lambda.
\]
Escribiremos entonces
\[
\log|\underline G|\leq\lambda.
\]

\subsubsection*{Lema 5.5.1.}
Sea $X$ una curva sobre $\mathbb F_q$, entendiendo aquí por ello un
esquema separado $X$ de tipo finito sobre $\mathbb F_q$ y de dimensión a
lo sumo $1$. Sea $\underline G$ un $E$-haz constructible sobre $X$ tal
que $\log|\underline G|\leq\lambda$, con $\lambda\in\mathbb R$. Entonces
los valores absolutos complejos de los valores propios $\alpha$ de $F$
sobre $\underline H_c^1(X,\underline G)$ verifican
\[
|\alpha|\leq q^{\lambda+1}.
\]

\pageoriginale

\noindent\underline{Demostración.}
Sea $\sigma$ una inmersión compleja de $E$, o simplemente una inmersión
compleja de la clausura algebraica de $\mathbb Q$ en $E$. Basta comprobar
que, para todo $\sigma$, los ceros $\beta$ de la función racional
$\sigma Z(X,\underline G,t)$ verifican
\[
|\beta^{-1}|\leq q^{\lambda+1}.
\]
En efecto, el Lema~5.2.1 muestra ya que, si $\alpha$ es un valor propio de
Frobenius sobre $\underline H_c^0(X,\underline G)$ o sobre
$\underline H_c^2(X,\underline G)$, entonces
$|\sigma\alpha|\leq q^{\lambda+1}$.

Comprobemos que la serie
\[
\log\sigma Z(X,\underline G,t)
=\sum_k\frac{t^k}{k}
\sum_{x\in X(\mathbb F_{q^k})}
\sigma\operatorname{Tr}(F,x^*\underline G)
\]
converge para $|t|<q^{-\lambda-1}$. Representando $X$ como un
recubrimiento ramificado de $\mathbb P^1$, se obtiene la estimación
elemental
\[
\#X(\mathbb F_{q^k})\leq Aq^k.
\]
Por otra parte, la hipótesis da
\[
\left|\sigma\operatorname{Tr}(F,x^*\underline G)\right|
\leq B(q^k)^\lambda.
\]
Así, el coeficiente $c_k$ de $t^k$ en
$\log\sigma Z(X,\underline G,t)$ verifica
\[
|c_k|\leq\frac1k C(q^{1+\lambda})^k,
\]
y se sigue la afirmación.

\subsubsection*{Teorema 5.5.2.}
Sea $X$ un esquema separado de tipo finito sobre $\mathbb F_q$, de
dimensión a lo sumo $n$, y sea $\underline G$ un $E$-haz constructible
sobre $X$ tal que $\log|\underline G|\leq\lambda$, con
$\lambda\in\mathbb R$. Entonces los valores absolutos complejos de los
valores propios $\alpha$ de Frobenius sobre
$\underline H_c^i(X,\underline G)$ verifican
\[
|\alpha|\leq q^{\lambda+i}
\qquad\text{y}\qquad
|\alpha|\leq q^{\lambda+n}.
\]

\pageoriginale

\noindent\underline{Demostración.}
Si $n\leq0$, el teorema es trivial. Si $n=1$, los casos $i=0$ e $i=2$
se deducen del Lema~5.2.1, mientras que el caso $i=1$ se deduce del
Lema~5.5.1. Se procede después por inducción sobre $n$, como en la
demostración del Teorema~5.2.2.

\subsubsection*{Corolario 5.5.3.}
Sea $X$ un esquema separado de tipo finito y dimensión $n$ sobre
$\mathbb F_q$. Los valores propios $\alpha$ del endomorfismo de Frobenius
geométrico $F$ sobre $\underline H_c^i(X,\mathbb Q_\ell)$ son enteros
algebraicos. Además:
\begin{enumerate}[label=(\roman*)]
\item si $\ell'\ne p$, $\alpha$ es una unidad $\ell'$-ádica;
\item si $i>n$, $\alpha$ es divisible por $q^{i-n}$;
\item si $i\leq n$ (respectivamente, $i\geq n$),
  $q^i\alpha^{-1}$ (respectivamente, $q^n\alpha^{-1}$) es un entero
  algebraico;
\item si $0<i$, los valores absolutos complejos de $\alpha$ verifican
  \[
  |\alpha|\leq q^{i-\frac12}
  \qquad\text{y}\qquad
  |\alpha|\leq q^n.
  \]
\end{enumerate}

\noindent\underline{Demostración.}
Por el Teorema~5.2.2, $\alpha$ y $q^{n-i}\alpha$ son enteros algebraicos.
Las afirmaciones (i)--(iii) se deducen del Teorema~5.4. Para $i>n$, (iv)
se deduce del Teorema~5.5.2. Para obtener (iv) cuando $i\leq n$, debemos
combinar el Teorema~5.5.2 con la hipótesis de Riemann para las curvas.

Procedemos por inducción sobre $n$, siendo trivial el caso $n=0$. Si $U$
es un abierto de $X$ cuyo complemento $Z$ tiene dimensión menor que $n$,
la sucesión exacta
\[
\underline H_c^i(U,\mathbb Q_\ell)\longrightarrow
\underline H_c^i(X,\mathbb Q_\ell)\longrightarrow
\underline H_c^i(Z,\mathbb Q_\ell)
\]
muestra que basta demostrar la afirmación para $U$.

\pageoriginale

Tras reemplazar primero $X$ por $X_{\mathrm{red}}$, nos reducimos al caso
en que $X$ es liso y afín y existe un morfismo liso
$f:X\longrightarrow S$ que verifica:
\begin{enumerate}[label=(\alph*)]
\item $\dim S\leq n-1$;
\item las fibras de $f$ son puramente unidimensionales.
\end{enumerate}
Como las fibras de $f$ son curvas afines,
\[
R^0f_!\mathbb Q_\ell=0.
\]
Por Weil [4],
\begin{equation}\tag{5.5.3.1}
\log\bigl|R^1f_!\mathbb Q_\ell\bigr|\leq\frac12,
\end{equation}
y el Lema~5.2.1 aplicado a las fibras de $f$ da
\begin{equation}\tag{5.5.3.2}
\log\bigl|R^2f_!\mathbb Q_\ell\bigr|\leq1.
\end{equation}
La sucesión espectral de Leray para $f$ se reduce aquí a una sucesión
exacta larga; su parte pertinente es
\[
\underline H_c^{i-2}(S,R^2f_!\mathbb Q_\ell)\longrightarrow
\underline H_c^i(X,\mathbb Q_\ell)\longrightarrow
\underline H_c^{i-1}(S,R^1f_!\mathbb Q_\ell).
\]
Por (5.5.3.1), respectivamente (5.5.3.2), y el Teorema~5.5.2, los valores
absolutos complejos de los valores propios $\alpha$ de Frobenius sobre el
término derecho, respectivamente el izquierdo, verifican
\[
|\alpha|\leq q^{i-1+\frac12}=q^{i-\frac12},
\]
respectivamente
\[
|\alpha|\leq q^{i-2+1}=q^{i-1}\leq q^{i-\frac12}.
\]
Se sigue el corolario.

\pageoriginale

\subsection*{5.6.}
Para obtener resultados análogos a los anteriores en cohomología sin
soportes parece necesario disponer de la resolución de singularidades.
En dimensión a lo sumo $1$, el problema queda resuelto por el
Corolario~5.3.

\subsubsection*{Teorema 5.6.}
Sea $N$ un entero, y supongamos disponible la resolución fuerte de
singularidades en el sentido de Hironaka para los esquemas de tipo finito
sobre $\mathbb F_q$ de dimensión a lo sumo $N$ (en la actualidad puede
tomarse, pues, $N=2$ [1]). Sea $f:X\longrightarrow Y$ un morfismo de
esquemas de tipo finito sobre $\mathbb F_q$, y sea $\underline G$ un
$E$-haz constructible $T$-entero sobre $X$. Si $\dim X\leq N$, entonces
los $E$-haces $R^if_*\underline G$ sobre $Y$ son $T$-enteros.

\noindent\underline{Demostración.}
Procedemos por inducción sobre la dimensión $n\leq N$ de $X$. Sea
$j:U\hookrightarrow X$ un abierto de $X$, liso y afín, tal que
$\underline G|_U$ sea constante torcido y el complemento $Z$ tenga
dimensión menor que $n$. Por resolución de singularidades, los haces
$R^ij_*(j^*\underline G)$ son constructibles (SGA~5, Exposé~I; [2]); si
$i>0$ están soportados sobre $Z$, y para $i=0$,
$j_*(j^*\underline G)$ solo difiere de $\underline G$ sobre $Z$. La
hipótesis de inducción muestra que basta demostrar el teorema para el haz
$j^*\underline G$ y los morfismos $j$ y $f\circ j$. Podemos suponer,
pues, que $X$ es liso y afín y $\underline G$ es constante torcido, es
decir, procede de una representación continua del grupo fundamental de
$X$ sobre un espacio vectorial de dimensión finita sobre $E$.

\pageoriginale

Sea $R$ el anillo de enteros de $E$, sea $\lambda$ un uniformizante, sea
$e$ el índice absoluto de ramificación, y elijamos un entero
$b>e/(\ell-1)$. Como $X$ es liso, luego normal, y su grupo fundamental es
profínito, existe un $R$-haz constante torcido $\underline G^\circ$ tal
que
\[
\underline G\simeq\underline G^\circ\otimes_R E.
\]
Sea $p:X'\longrightarrow X$ el recubrimiento étale que trivializa el haz
localmente constante $\underline G^\circ\otimes_RR/\lambda^b$. El haz
$\underline G$ es un factor directo de $p_*p^*\underline G$, de modo que
basta demostrar el teorema para $p^*\underline G$ y $f\circ p$. Queda
tratar el caso en que se cumplen las condiciones siguientes:
\[
(*)\quad
\begin{minipage}{0.82\linewidth}
$X$ es liso y afín, y existe un $R$-haz constante torcido
$\underline G^\circ$ tal que
$\underline G\simeq\underline G^\circ\otimes_R E$ y
$\underline G^\circ\otimes_RR/\lambda^b$ es constante.
\end{minipage}
\]

Factoricemos $f$ como un morfismo propio seguido de una inmersión abierta,
$f=\bar f\circ j$:
\begin{center}
\begin{tikzcd}
X \arrow[r,hook,"j"'] \arrow[d,"f"'] & X' \arrow[dl,"\bar f"] \\
Y & {}
\end{tikzcd}
\qquad .
\end{center}
El teorema ya se ha demostrado para $\bar f$ (5.2), por lo que basta
demostrarlo para $j$. Aplicando la resolución de singularidades a $X'$,
podemos suponer que $X'$ es liso y que $X$ es el complemento en $X'$ de
un divisor con cruces normales, unión de divisores lisos. La hipótesis
$(*)$ implica que $\underline G$ está moderadamente ramificado a lo largo
de este divisor. Tenemos entonces:

\pageoriginale

\subsubsection*{Lema 5.6.1.}
Sea $X$ un esquema liso de tipo finito sobre $\mathbb F_q$, sea $D$ un
divisor con cruces normales en $X$, unión de divisores lisos $D_i$, sea
$j:U\hookrightarrow X$ el complemento de $D$ en $X$, y sea $\underline G$
un $E$-haz constante torcido sobre $U$, moderadamente ramificado a lo
largo de $D$. Si $\underline G$ es $T$-entero, entonces los haces
$R^ij_*\underline G$ son $T$-enteros.

Admitiremos el resultado siguiente (cf.~SGA~1, Exposé~XIII).

\subsubsection*{Lema 5.6.2.}
\begin{enumerate}[label=(\roman*)]
\item Bajo las hipótesis del Lema~5.6.1, sea $D'$ una componente de $D$,
sea $D''$ la unión de las demás componentes, y pongamos
$U'=X-D''$ y $D^{\prime\circ}=D'\cap U'$:
\begin{center}
\begin{tikzcd}
{} & D^{\prime\circ}
 \arrow[r,hook,"j'_D"'] \arrow[d,"i^\circ"] &
 D' \arrow[d,"i"] \\
U \arrow[r,hook,"j_1"'] &
 U' \arrow[r,hook,"j'"'] & X
\end{tikzcd}
\qquad .
\end{center}
Entonces la aplicación de cambio de base
\[
i^*Rj'_*(Rj_{1*}\underline G)
\longrightarrow
Rj'_{D*}\,i^{\circ *}(Rj_{1*}\underline G)
\]
es un isomorfismo, y los haces
$i^{\circ *}Rj_{1*}\underline G$ son constantes torcidos y moderadamente
ramificados a lo largo de $D'-D^{\prime\circ}$.

\item Si $C$ es una curva lisa sobre $X$, transversal a $D$ en un punto
liso $x$ de $D$,
\begin{center}
\begin{tikzcd}
C\cap U \arrow[r,"j_c"'] \arrow[d,"k'"] & C \arrow[d,"k"] \\
U \arrow[r,"j"'] & X
\end{tikzcd}
\qquad ,
\end{center}
entonces la aplicación de cambio de base
\[
k^*Rj_*\underline G\longrightarrow Rj_{c*}\,k^{\prime *}\underline G
\]
es un isomorfismo en $x$.
\end{enumerate}

Deducimos el Lema~5.6.1 del Lema~5.6.2 por inducción sobre la dimensión,
admitiendo provisionalmente el Lema~5.6.1 para las curvas. El
Lema~5.6.2(ii) y esta hipótesis muestran que los haces
$R^ij_*\underline G$ son $T$-enteros en todo punto liso de $D$. En
particular, en la notación del Lema~5.6.2(i), los haces
$R^ij_{1*}\underline G$ son $T$-enteros.

\pageoriginale

El Lema~5.6.2(i), junto con la hipótesis de inducción aplicada a
$i^{\circ *}Rj_{1*}\underline G$ y a $j'_D$, muestra entonces que los
haces $i^*R^ij_*\underline G$ son $T$-enteros, cualquiera que sea la
componente $D'$ elegida. Esto demuestra el Lema~5.6.1.

Queda tratar las curvas. Bajo las hipótesis del Lema~5.6.1, sea $x$ un
punto de $X-U$, sea $\mathfrak G$ el grupo de Galois del cuerpo de
fracciones de $\widehat{\mathcal O}_{x,X}$, y sea $\mathfrak J$ el grupo
de inercia. Denotemos también por $\underline G$ la representación de
$\mathfrak G$ definida por el haz, y sea $\bar x$ un punto geométrico de
$X$ localizado en $x$. Hay isomorfismos de
$\mathfrak G/\mathfrak J$-módulos galoisianos
\[
(j_*\underline G)_{\bar x}\simeq\underline G^{\mathfrak J},
\]
\[
(R^1j_*\underline G)_{\bar x}\simeq\underline G_{\mathfrak J}(-1),
\]
y
\[
(R^ij_*\underline G)_{\bar x}=0\qquad\text{para }i\geq2.
\]
La afirmación se deduce del Corolario~5.3(a) y del Lema~5.3.1.

\pageoriginale

\section*{Bibliografía}
\begin{enumerate}[label={\textup{[\arabic*]}},leftmargin=2.4em]
\item S.~S. Abhyankar, \emph{Resolution of Singularities of Embedded
  Algebraic Surfaces}, Academic Press, Nueva York, 1966.
\item J.-P. Jouanolou, \emph{Cohomologie $\ell$-adique des schémas}, de
  próxima aparición en \emph{Lecture Notes in Mathematics}.
\item S. Lang y A. Weil, «Number of Points of Varieties in Finite
  Fields», \emph{Amer. J. Math.} \textbf{76} (1954), 819--827.
\item A. Weil, \emph{Sur les courbes algébriques et les variétés qui
  s'en déduisent}, \emph{Actualités Sci. Indust.} n.º~1041, Hermann,
  París, 1948.
\end{enumerate}
